\begin{textsample}{POS Dim 2 – vem\_ed\_21 – Score 5.00 – vem\_ed\_21/t108.txt}  \label{ex:f2_pos_168}
continuação Apoio ao \textbf{matchmaking} entre parceiros, \textbf{treinamento} docente, orientação e \textbf{suporte}, reconhecimento de professores e alunos estão entre as recomendações para \textbf{institucionalizar} as iniciativas COIL.
A segunda comunicação de Ana Carolina foi"Regional campuses in the lead: Institutionalizing COIL across multi-campus public systems"(Campi regionais na liderança: \textbf{institucionalizando} COIL em sistemas públicos multicampi), com Dan Nolan (University of Minnesota System / EUA) e Natalia Dyba (University of Washington Bothell / EUA).
Os autores relataram os esforços de institucionalização de campi regionais em sistemas públicos no Brasil e nos Estados Unidos.
Destacaram que a organização dos Intercâmbios Virtuais nesse contexto requer alinhamento com metas \textbf{interinstitucionais} mais amplas e uma noção dos pontos fortes, da preparação e da capacidade do corpo docente, da equipe e da liderança.

% matched lemmas: institucionalizar, interinstitucional, matchmaking, suporte, treinamento
\end{textsample}
